\subsection*{C.2 \quad What the rotation costs}

Without the three rules the discounted margin is $357\,029.70$ EUR. \textbf{The rules cost $17\,096$ EUR, i.e.\ $4.79\%$.}

\begin{figure}[!ht]
\centering
\begin{tikzpicture}
\begin{axis}[width=10cm, height=5cm, xbar, bar width=8pt, symbolic y coords={only legume,only preceding crop,only return period,without legume,without preceding crop,without return period,all three rules,no rules}, ytick=data, xmin=0, xmax=21000, scaled x ticks=false, xticklabel style={/pgf/number format/.cd,fixed,precision=0,1000 sep={}}, xlabel={loss against ``no rules'' (EUR)}, nodes near coords, nodes near coords style={font=\tiny,/pgf/number format/.cd,fixed,precision=0,1000 sep={}}, enlarge y limits=0.1]
\addplot[fill=cbeet] coordinates {(9711.06,only legume)(1479.86,only preceding crop)(9106.76,only return period)(9106.76,without legume)(17096.04,without preceding crop)(10618.24,without return period)(17096.04,all three rules)(0,no rules)};
\end{axis}
\end{tikzpicture}
\caption{Cost of the rules, in each case against the plan without any rule ($357\,029.70$ EUR).}
\end{figure}

\begin{table}[!ht]
\centering\footnotesize
\begin{tabular}{lrrr}
\toprule
case & discounted margin (EUR) & cost (EUR) & cost (\%)\\
\midrule
all three rules & 339\,933.66 & 17\,096.04 & 4.79\\
without the return period & 346\,411.47 & 10\,618.24 & 2.97\\
without the preceding crop & 339\,933.66 & 17\,096.04 & 4.79\\
without the legume rule & 347\,922.94 & 9\,106.76 & 2.55\\
only return period & 347\,922.94 & 9\,106.76 & 2.55\\
only preceding crop & 355\,549.84 & 1\,479.86 & 0.41\\
only legume & 347\,318.64 & 9\,711.06 & 2.72\\
no rule & 357\,029.70 & 0 & 0\\
\bottomrule
\end{tabular}
\caption{Discounted margin in each case; the cost is the difference to ``no rule''.}
\end{table}


\textbf{Which rule is expensive?} Remove one rule at a time:
\begin{itemize}
\item without legumes the farm gains $7\,989$ EUR;
\item without the return period it gains $6\,478$ EUR;
\item without the preceding-crop rule it gains nothing.
\end{itemize}
The legume rule is the expensive one. It costs more by both measures: $7\,989$ against $6\,478$ EUR when removed alone, and $9\,711$ against $9\,107$ EUR when kept alone.

\textbf{Why the plan resists:}
\begin{itemize}
\item \textbf{Return period.} The free plan repeats each crop on its best soil every year: $14$ ha of wheat on B4 (heavy clay), $14.9$ ha of potato and $11.1$ ha of beet on B1 (deep loam). The rule caps a crop over two consecutive years at $r_cA_b$, e.g.\ $12.6$ ha of wheat on B4 and $13$ ha of potato on B1, so only about half of these areas can stay on their best block each year. The rest moves to blocks with lower yields: potatoes to B2 and B3, barley and beans onto B4.
\item \textbf{Legume.} The free plan grows no beans: they have the lowest margin on every soil ($351$--$554$ EUR/ha, against at least $661$ EUR/ha for any other productive crop). The rule forces $8$ ha of beans every year in place of more profitable crops.
\item \textbf{Preceding crop.} With the other two rules it is not binding: the return period already moves potatoes from B1 to B2 and B3, which carried wheat the year before. Alone it costs $1\,480$ EUR: the free plan grows potatoes on B1 every year with no cereal before them, so the farm must put wheat on B1 in some years to keep its potatoes there.
\end{itemize}
